\section{Related Work}
\label{sec:related_work}

\subsection{Vision-Language-Action Models}

Recent advances in foundation models have enabled a new class of robotic policies that map visual observations and language instructions directly to actions. RT-2~\cite{brohan2023rt2} first demonstrated that fine-tuning large vision-language models on robot demonstrations transfers web-scale semantic knowledge to robotic control; it tokenizes actions as text and trains jointly on internet and robot data. Octo~\cite{ghosh2024octo} introduced a modular transformer backbone trained on 800k trajectories from the Open X-Embodiment dataset, with flexible adaptation to new observation spaces and embodiments. OpenVLA~\cite{kim2024openvla} scaled this paradigm to 7B parameters with a fused SigLIP-DinoV2 visual encoder and an MLP projector that maps visual embeddings into the Llama~2 input space; it outperforms RT-2-X (55B) by 16.5\% in absolute success rate with 7$\times$ fewer parameters. SmolVLA~\cite{shukor2025smolvla} pursued the opposite end of the efficiency spectrum: a 450M-parameter model that combines aggressive pixel shuffle compression (64 visual tokens per frame) with a flow-matching action expert, and matches or exceeds larger VLAs on both simulation and real-world benchmarks. $\pi_0$~\cite{black2024pi0} and $\pi_{0.5}$~\cite{physicalintelligence2025pi05} built on PaLiGemma with a separate action expert connected via cross-attention and produce smooth 50Hz action trajectories through flow matching.

A critical but underexplored dimension across these models is the \emph{vision connector}: the module that bridges the visual encoder and the language/action backbone. OpenVLA and LLaVA~\cite{liu2023llava} use per-patch MLP projectors that preserve the spatial independence of each visual token. SmolVLA and InternVL~\cite{chen2024internvl} employ pixel shuffle (space-to-depth) connectors that fold spatial neighborhoods into channel dimensions before projection, creating tokens whose information content spans multiple spatial locations. BLIP-2~\cite{li2023blip2} introduced the Q-Former, a learned cross-attention bottleneck with fixed-size query vectors, while Flamingo~\cite{alayrac2022flamingo} proposed the Perceiver Resampler for a similar purpose. These architectural choices determine how visual information is distributed across tokens, yet no prior work has systematically studied how connector type affects the compressibility of the resulting token sequences in the context of action generation.

\subsection{Token Efficiency in Vision-Language Models}

Reducing the computational cost of visual tokens has been extensively studied in the vision transformer and VLM literature. DynamicViT~\cite{rao2021dynamicvit} introduced learnable prediction modules that progressively prune tokens at multiple transformer layers and reduce FLOPs by 30\% with less than 0.5\% accuracy loss. Token Merging (ToMe)~\cite{bolya2023tome} takes a complementary approach, bipartite-matching similar tokens and merging them rather than discarding, which preserves foreground information that pruning-based methods may lose. EViT~\cite{liang2022evit} reorganizes tokens by fusing inattentive patches into a single representative token, guided by class token attention. FastV~\cite{chen2024fastv} extended these ideas to multimodal LLMs and showed that 50\% of visual tokens can be pruned after shallow layers with negligible performance loss on image and video understanding tasks.

In the VLA setting, several concurrent works have adapted token pruning for robotic control. VLA-IAP~\cite{cheng2026vlaiap} proposed an interaction-aligned pruning strategy that preserves structurally critical manipulation regions through a geometric prior and retains 97.8\% success rate on OpenVLA with 1.25$\times$ speedup. VLA-Pruner~\cite{liu2025vlapruner} introduced dual-level importance scoring that combines semantic prefill attention with action decode attention estimated via temporal smoothing. VLA-Cache~\cite{xu2025vlacache} exploits temporal coherence between consecutive frames to reuse KV-cache entries for unchanged visual tokens. At the LLM level, H$_2$O~\cite{zhang2023h2o} identified ``heavy hitter'' tokens that dominate attention scores and proposed a KV-cache eviction policy that retains only these high-utility entries.

A common limitation of these methods is that they were developed and evaluated on a single architecture (typically OpenVLA or a standard VLM with MLP projection). Whether the same pruning strategies transfer across architectures with fundamentally different vision connectors remains untested. Our work fills this gap by showing that connector choice determines whether token pruning produces graceful degradation or catastrophic failure.

\subsection{Vision Connectors and Information Density}

The vision connector governs how spatial information from the visual encoder is reorganized before entering the language model, and this reorganization has direct consequences for per-token information density. MLP projectors, as used in LLaVA~\cite{liu2023llava} and OpenVLA~\cite{kim2024openvla}, apply independent linear transformations to each patch embedding. Because each output token corresponds to exactly one spatial location, removing a token discards information from that location alone, and the remaining tokens retain their original semantics. This locality property makes MLP-projected tokens amenable to importance-based pruning: low-attention tokens can be dropped with bounded information loss.

Pixel shuffle connectors, adopted by SmolVLA~\cite{shukor2025smolvla} and InternVL~\cite{chen2024internvl}, operate differently. The space-to-depth rearrangement groups $k \times k$ spatial neighborhoods and concatenates their features along the channel dimension before a linear projection. The resulting tokens each encode a mixture of $k^2$ original patch representations. This creates two related problems for token pruning. First, tokens are no longer spatially independent: dropping one token removes information from multiple spatial positions, and the lost information cannot be recovered from neighboring tokens because they encode \emph{different} spatial neighborhoods. Second, the information content per token is higher and more uniformly distributed, so the assumption that some tokens carry negligible information (the basis for attention-based pruning) breaks down. We refer to this phenomenon as \emph{token entanglement} and show empirically that it produces a cliff-edge failure mode: even modest pruning ratios cause complete task failure (0\% success rate) on SmolVLA, in contrast to the graded degradation observed on OpenVLA.

Learned compression modules such as the Q-Former~\cite{li2023blip2} and Perceiver Resampler~\cite{alayrac2022flamingo} represent a third paradigm. These modules use a fixed set of learnable queries that cross-attend to the full set of visual features, and produce a compact representation trained end-to-end to preserve task-relevant information. Unlike pixel shuffle, the compression is adaptive: the cross-attention weights determine which spatial regions contribute to each output token based on learned relevance rather than fixed spatial grouping. This adaptivity makes learned compression a natural candidate for addressing the token entanglement problem in pixel shuffle architectures.

\subsection{Learned Token Compression}

Cross-attention-based compression was introduced in the vision-language literature to reduce the number of visual tokens entering the language model. The Q-Former in BLIP-2~\cite{li2023blip2} uses 32 learnable queries in a 12-layer transformer that cross-attends to frozen image encoder outputs; it acts as an information bottleneck trained in two stages, first for vision-language representation learning and then for generative pre-training. The Perceiver Resampler in Flamingo~\cite{alayrac2022flamingo} serves a similar function: it maps variable-length spatio-temporal visual features to a fixed number of output tokens that condition the language model through gated cross-attention layers.

These architectures were designed to bridge modality gaps in VLMs, not to address the specific challenges of compressing entangled token representations in VLA models. Our Smart Compress module draws on the same cross-attention mechanism but applies it to a different problem: recovering compressibility in architectures where post-connector tokens resist importance-based pruning. By training a small set of queries to cross-attend to the pixel-shuffle output, the module learns to disentangle the spatial information that was mixed during the space-to-depth rearrangement and yields tokens that can be pruned without catastrophic failure. This positions learned compression not as an alternative to token pruning, but as a prerequisite for making pruning viable on entangled architectures.
